% Resumo em língua estrangeira
\begin{resumo}[Abstract]
\begin{otherlanguage*}{english}
Simply supported steel-concrete composite beams with profiled steel deck or solid slabs, built without shoring, are the prevailing floor solution in Brazilian commercial buildings. Their design is multi-phase: during casting (Phase 1) the bare steel beam carries its self-weight, the wet concrete and the construction load; after hardening (Phase 2) the section acts compositely, with longitudinal shear transferred by connectors. The verification covers ultimate limit states (bending, shear and degree of shear connection) and serviceability limit states (deflections including creep and shrinkage and, since the 2024 edition of ABNT NBR 8800, walking-induced floor vibration in Annex I). This work develops an open-source computational tool, written in C++20 with a Dear ImGui and Raylib graphical interface, that automates this chain for both the 2008 and the 2024 editions of the code, keeping the two implementations separate so that they can be compared directly. The theoretical background of Annex L, the formalisation of the algorithms and the software verification strategy based on automated tests are presented. The tool was validated against an office-floor example (W460\ensuremath{\times}52.0 beam, MF-75 steel deck, 9.0 m span) solved in a reference spreadsheet: the 25 compared quantities (internal forces, resistances, degree of connection, effective inertias, deflections and vibration parameters) agreed within a maximum relative difference of $2.4 \times 10^{-15}$, and the vibration procedure was also checked against a second, independent example. Parametric studies showed that, for the case analysed, the 2024 edition lowers the minimum degree of shear connection from 0.519 to 0.280 but, for solid slabs, requires about 30\% more connectors for full interaction because of the new $R_p = 0.75$ factor; they also showed that the indirect vibration criterion of the 2008 edition rejects the case-study beam, while the Annex I assessment accepts it under the peak-acceleration criterion. The modular architecture lays the groundwork for a future extension to prestressed composite beams.

\textbf{Keywords:} Steel-concrete composite beams; ABNT NBR 8800:2024; Computational tool; Unshored construction; Steel deck; Degree of shear connection; Floor vibration.
\end{otherlanguage*}
\end{resumo}
